\subsection*{Notes for III. POSITIONING IN LITERATURE}
\paragraph{Objective}
Explain where P3Net fits relative to (i) multiobjective NAS (NSGA-Net/NSGANet), (ii) HPO/BO methods in discrete spaces (TPE/SMAC/BOHB), (iii) EDAs / probabilistic model-building EAs (P3), and (iv) multi-fidelity NAS/HPO.

\paragraph{Keep}
\begin{itemize}
  \item The intent is to synthesize and justify novelty/comparability.
\end{itemize}

\paragraph{TODO:}
\begin{itemize}
    \item What is missing (add)
\begin{itemize}
  \item Multiobjective NAS lineage: what NSGANet/NSGANet does (encoding, selection/crowding, operators, objectives).
  \item Background of the P3 algorithm: what distribution family does it represent, how the dependency structure is learned (MST), and how sampling works.
  \item EDAs for structured/discrete optimization: why dependency modeling matters (vs independent categorical sampling).
  \item Multi-fidelity NAS/HPO: ASHA/Hyperband and why mixing fidelities can bias learning/selection.
\end{itemize}
\end{itemize}
\begin{itemize}
    \item \paragraph{Add first}
\end{itemize}
\begin{itemize}
  \begin{itemize}
      \item A short taxonomy (3--5 paragraphs) followed by a “gap + thesis” paragraph that motivates P3Net.
      \item End with “What exactly is new in P3Net” in bullet form (integration + fidelity-consistent learning + NAS validity/diversity).
  \end{itemize}
\end{itemize}

\paragraph{Peer-review must-haves}
\begin{itemize}
  \item A rationale for a defensible baseline set (why each is relevant, feasible, and fairly adapted).
\end{itemize}
